\akapit Stanowisko symulatora bezzałogowego statku powietrznego składa się z systemu wizualizacji, kontrolerów lotu oraz panelu instruktora[cite: 4]. System wizualizacji obejmuje dwa duże monitory prezentujące środowisko symulacyjne z perspektywy pierwszej osoby (kamera na pokładzie) oraz z zewnątrz (widok pogoni), a także podstawowe wskaźniki przyrządowe (prędkościomierz, sztuczny horyzont, wysokościomierz, chyłomierz, wskaźnik kursu i wariometr)[cite: 4]. Ręczne sterowanie modelem realizowane jest za pomocą zestawu sterowniczego składającego się z joysticka i przepustnicy[cite: 4].

\akapit Zarządzanie symulacją, rejestracja danych oraz zmiana parametrów odbywają się z poziomu panelu instruktora, wyświetlanego na jednym z mniejszych monitorów i obsługiwanego za pomocą myszy oraz klawiatury[cite: 4, 5]. Pod względem matematycznym model symulacyjny opiera się na nieliniowych równaniach ruchu zdefiniowanych w trzech układach współrzędnych: inercyjnym ziemskim ($O_n x_n y_n z_n$), grawitacyjnym ($O_g x_g y_g z_g$) oraz związanym z obiektem ($O_b x_b y_b z_b$)[cite: 6, 7]. W strukturze funkcjonalnej modelu wyróżnia się blok dynamiki samolotu, model układu sterowania (uwzględniający dynamikę serwomechanizmów, luzy i opóźnienia), model systemu pomiarowego (szumy czujników) oraz blok autopilota odpowiedzialnego za sterowanie po zadanej trajektorii[cite: 7, 8].
